%=============================================================================!
% 11.8  WHAT THE TRAJECTORY LOOKS LIKE                                        !
%=============================================================================!
\section{What the trajectory looks like}
\label{sec:cs-trajectory}

A run produces a record: energies at every step, positions and
velocities every so often. Before any of it is used, the record should be
read for signs that the run went wrong, and many faults show in it
clearly when the reader knows where to look. This section shows a
healthy run of rigid water and the same run broken in eight ways, each
caught by a check of \texttt{mdlab.diagnostics} or of the record itself. Later chapters add the
checks of temperature, pressure and thermostats, and Chapter~17 collects
them all.

\subsection*{A healthy run}

\Cref{fig:cs-healthy} follows the 64 rigid molecules for
\SI{10}{\pico\second}, 5000 steps of \SI{2}{\femto\second} by RATTLE. The
kinetic and potential energies swing by up to about
\SI{13}{\milli\electronvolt} per molecule over the first
\SI{2}{\pico\second}, in opposite directions, so that their sum, the total
energy, stays in a band \SI{0.31}{\milli\electronvolt} wide per molecule
over the whole run, \SI{0.020}{\electronvolt} for the box: the mirror image
of $K$ and $U$ is the first sign of a healthy run without a thermostat.
The standard deviation of the total energy is 0.93\% of the kinetic
energy's, just inside the 1\% level chosen in \cref{sec:cs-water}. The
mean total energy over the second half of the run, less that over the
first, is \SI{1.1e-5}{\electronvolt}, and divided by the
\SI{5}{\pico\second} between the centres of the halves it gives a drift of
\SI{2.2e-6}{\electronvolt\per\pico\second}, far below the fluctuation;
this is what \texttt{energy\_drift} computes. Every held distance stays
within \SI{1.5e-10}{\angstrom} of its length, the tolerance of $10^{-10}$
times the H--H length, and the centre of mass moves
\SI{1.8e-10}{\angstrom} in \SI{10}{\pico\second}, the rounding of the
arithmetic. The root-mean-square displacement of the oxygen atoms, the
square root of the mean over them of the squared distance each has moved
from its start, grows steadily, \num{1.41}, 4.14 and
\SI{5.76}{\angstrom} after 1, 5 and \SI{10}{\pico\second}, the wandering
of molecules in a liquid that Chapter~16 turns into a diffusion
coefficient. Over the run the kinetic energy averages
\SI{4.99}{\electronvolt}, \SI{13.10}{\milli\electronvolt} per freedom,
with a standard deviation of \SI{0.29}{\electronvolt}; the potential
energy averages \SI{-0.403}{\electronvolt} per molecule. These are
averages over one run of \SI{10}{\pico\second}, given with their spread;
how far such averages can be trusted, their error bars, is the subject of
Chapter~15. \Cref{fig:cs-frames} shows three
frames, with three molecules marked to follow them.

\begin{figure}[htb]
\centering
\includegraphics{ch11_constraints/healthy}
\caption{A healthy run: 64 rigid water molecules followed by RATTLE for
\SI{10}{\pico\second} with $\dt = \SI{2}{\femto\second}$. (a) The changes
of the kinetic, potential and total energies from their starting values,
per molecule, over the first \SI{2}{\pico\second}. (b) The largest error
of a held distance and the distance the centre of mass has moved, on a
logarithmic axis. (c) The root-mean-square distance the oxygen atoms have
moved from their starts.}
\label{fig:cs-healthy}
\end{figure}

\begin{figure}[htb]
\centering
\includegraphics{ch11_constraints/frames}
\caption{Frames of the healthy run at 0, 5 and \SI{10}{\pico\second},
each molecule moved whole into the cell by its oxygen and drawn with its
two O--H bonds (oxygen red, hydrogen white). Three molecules have their
oxygen drawn in teal so that their motion can be followed. All frames are
seen from the same direction at the same scale.}
\label{fig:cs-frames}
\end{figure}

\subsection*{The same run, broken}

\Cref{fig:cs-faults} breaks the run in eight ways, each a mistake that is
easy to make, and \cref{tab:cs-faults} lists the check that catches each,
with its value in the healthy and the broken run.

\begin{figure}[p]
\centering
\includegraphics{ch11_constraints/faults}
\caption{The run of \cref{fig:cs-healthy} broken in eight ways (teal),
beside the healthy run (grey dashed), each through the quantity that
catches the fault. (a) Flexible water at \SI{2}{\femto\second} against
\SI{0.5}{\femto\second}. (b) The Coulomb energy cut off atom by atom at
\SI{6}{\angstrom} against the Ewald sum. (c) Velocities in
\si{\angstrom\per\femto\second} read as \aa ngstr\"oms per ASE time unit:
the kinetic energy per freedom against the value asked for (dotted).
(d) Positions wrapped into the cell at every step while the bonds are
measured without their nearest copies: the largest bond error. (e) One
molecule placed with its oxygen \SI{1}{\angstrom} from another's: the
forces at the start, largest first. (f) Constraints held to $10^{-3}$
against $10^{-10}$. (g) Starting velocities with their total momentum
left in. (h) RESPA with outer steps of \SI{4}{\femto\second} against
\SI{1}{\femto\second}.}
\label{fig:cs-faults}
\end{figure}

\begin{table}[htb]
\centering
\caption{The faults of \cref{fig:cs-faults}, the check of
\texttt{mdlab.diagnostics} or of the record that catches each, and its
value in the healthy run and in the broken one; runs of
\SI{2}{\pico\second}, except the flexible run at \SI{0.5}{\femto\second}
and both RESPA runs, of \SI{1}{\pico\second}, and runs that failed
sooner.}
\label{tab:cs-faults}
\small
\begin{tabular}{@{}llll@{}}
\toprule
fault & check & healthy & broken\\
\midrule
(a) step too long & \texttt{energy\_fluctuation} & 1.5\% & 47\%\\
(b) jumping cutoff & \texttt{energy\_fluctuation} & 1.0\% & 2100\%\\
(c) wrong units & $K(0)$ per freedom &
   \SI{12.9}{\milli\electronvolt} & \SI{0.125}{\milli\electronvolt}\\
(d) wrapped wrongly & \texttt{bond\_error} &
   \SI{1.5e-10}{\angstrom} & \SI{15.8}{\angstrom}\\
(e) overlap & \texttt{largest\_force} &
   \SI{2.9}{\electronvolt\per\angstrom} &
   \SI{3.0e5}{\electronvolt\per\angstrom}\\
(f) loose constraints & \texttt{energy\_drift} &
   \SI{4e-4}{\electronvolt\per\pico\second} &
   \SI{-0.011}{\electronvolt\per\pico\second}\\
(g) drift kept & \texttt{centre\_of\_mass\_path} &
   \SI{8e-12}{\angstrom} & \SI{0.99}{\angstrom}\\
(h) resonance & \texttt{energy\_drift} &
   \SI{2e-4}{\electronvolt\per\pico\second} &
   \SI{2.2}{\electronvolt\per\pico\second}\\
\bottomrule
\end{tabular}
\end{table}

\emph{(a) Too long a step.} Flexible water run at \SI{2}{\femto\second},
the step that suits rigid water, has an energy band 47\% of the kinetic
energy's spread, against 1.5\% at \SI{0.5}{\femto\second}, and a mean that
moves by \SI{0.016}{\electronvolt\per\pico\second}. The check is the
energy fluctuation, and the step test of \cref{sec:in-stability} confirms
the cause: halving the step should cut the fluctuation by about four, but
from 2 to \SI{1}{\femto\second} it cuts it by 7.5, so
\SI{2}{\femto\second} lies beyond the range where it grows as $\dt^2$.

\emph{(b) A cutoff that jumps.} Replacing the Ewald sum by the Coulomb
energy of pairs closer than \SI{6}{\angstrom}, cut off atom by atom, makes
the energy jump whenever a pair crosses the cutoff: an oxygen and a
hydrogen of different molecules have
$k_{\mathrm e}q_{\mathrm O}q_{\mathrm H}/r = \SI{-0.835}{\electronvolt}$
there. The total energy leaps by \SI{48}{\electronvolt} in the first
\SI{18}{\femto\second} and then wanders over some \SI{130}{\electronvolt},
its standard deviation 21 times the kinetic energy's, while its mean moves
by \SI{-8.2}{\electronvolt\per\pico\second}: the fault of
\cref{sec:pe-cutoffs} at its worst. A band or a drift that does not shrink
as the step is shortened points to forces that are not the slopes of a
smooth energy.

\emph{(c) Wrong units.} ASE measures time in its own unit,
\aa ngstr\"oms times $\sqrt{\mathrm{amu}/\mathrm{eV}}$, which is
$1/\sqrt{\texttt{FORCE\_TO\_ACCEL}} = \SI{10.18}{\femto\second}$ with the
factor of \cref{sec:ne-atoms}. Velocities written in
\si{\angstrom\per\femto\second} and read by a program that takes them in
that unit are 10.18 times too small, and the kinetic energy $10.18^2 =
104$ times too small: \SI{0.125}{\milli\electronvolt} per freedom at the
start against the \SI{12.93}{\milli\electronvolt} asked for. Nothing else
looks wrong: the energy is conserved as well as in the healthy run, with a
fluctuation of 0.6\%, and the kinetic energy rises within a tenth of a
picosecond as the warm arrangement gives up energy, settling near
\SI{8.2}{\milli\electronvolt} per freedom. Conservation shows only that a
run is consistent with itself; whether it is the run asked for is shown by
comparing the kinetic energy at the start with the value requested.

\emph{(d) Wrapped where it must not be.} Wrapping each atom into the cell
at every step is harmless for the forces, which take nearest copies, but
not for a code that measures a bond as $\vb{r}_j - \vb{r}_i$ directly: a
molecule that lies across a face of the cell then seems to have a bond as
long as the cell, or longer if it lies across an edge. Here the prepared
positions spread over 2.2 cell widths, so the first wrap splits
molecules, the largest bond error jumps from $\num{1.5e-10}$ to
\SI{15.8}{\angstrom} at the first step, and the run fails at its third.
The opposite fault, positions left outside the cell where a code needs
them inside it, is why the Ewald sum of \cref{sec:md-ewald} wraps them
first. Displacements, by contrast, must be measured from positions that
have not been wrapped: \texttt{largest\_jump} reports the largest move
between saved frames as a fraction of the cell's width, 0.16 in the
unwrapped record of the healthy run, and near 1 when an atom has been
wrapped between them.

\emph{(e) Overlapping atoms at the start.} One molecule placed with its
oxygen \SI{1}{\angstrom} from another's brings two atoms
\SI{0.886}{\angstrom} apart, against \SI{1.52}{\angstrom} for the closest
atoms of different molecules in the healthy start, a potential energy of
\SI{+25\,192}{\electronvolt} and a largest force of
\SI{3.0e5}{\electronvolt\per\angstrom} against \SI{2.9}{\electronvolt
\per\angstrom}. The first step fails. The checks belong before the run:
\texttt{closest\_approach}, leaving out bonded pairs, and
\texttt{largest\_force} on the starting frame.

\emph{(f) Constraints drifting.} Stopping SHAKE and RATTLE at a tolerance
of $10^{-3}$ leaves bond errors of \SI{1.5e-3}{\angstrom}; the energy falls
by \SI{0.093}{\electronvolt} in the first step, as every molecule is
distorted a little, and then drifts by
\SI{-0.011}{\electronvolt\per\pico\second}, against
\SI{4e-4}{\electronvolt\per\pico\second} in the healthy run over the same
\SI{2}{\pico\second}. The energy gives the alarm, and the largest bond
error, \SI{1.5e-3}{\angstrom} against \SI{1.5e-10}{\angstrom}, the cause: a
bond error at the tolerance asked for is healthy only if a tighter
tolerance leaves the energy unchanged.

\emph{(g) Drift of the centre of mass.} Starting velocities drawn without
removing their total momentum (\cref{sec:md-start}) leave the whole box
moving: here $\vb{P} = (0.28, 0.29, 0.40)$\,amu\,\AA/fs, so the centre of
mass moves at $\lvert\vb{P}\rvert/M = \SI{0.50}{\angstrom\per\pico\second}$
and has gone \SI{0.99}{\angstrom} in \SI{2}{\pico\second}, against
\SI{8e-12}{\angstrom} in the healthy run. The energy gives no sign, since
momentum is conserved and its kinetic energy, \SI{0.015}{\electronvolt},
never changes; that kinetic energy belongs to no motion of the molecules
relative to one another. \texttt{centre\_of\_mass\_path} catches it.

\emph{(h) Resonance.} RESPA with outer steps of \SI{4}{\femto\second},
near half the period of the model's stretches, gains
\SI{2.2}{\electronvolt\per\pico\second}, against
\SI{2e-4}{\electronvolt\per\pico\second} with outer steps of
\SI{1}{\femto\second}, and its band is 2.5 times the kinetic energy's
spread (\cref{sec:cs-resonance}).

\begin{practice}
A run is read in three passes. Before it starts: the closest approach and
the largest force on the starting frame, the kinetic energy against the
value asked for, and the total momentum. While it runs: the total energy,
whose fluctuation should be small against the kinetic energy's and fall
by about four when the step is halved, and whose mean should not move;
the bond errors, at a tolerance tight enough that tightening it leaves
the energy unchanged; and the centre of mass, still. After it: the
positions used for displacements must be continuous, with every move
between frames a small fraction of the cell. These checks find faults;
whether the run samples what it should is the subject of Chapters~12
to~15.
\end{practice}

\begin{keyidea}
In a healthy run without a thermostat, $K$ and $U$ mirror each other and
their sum stays in a narrow band with no drift, held distances stay at
the tolerance and the centre of mass stays still. Each common fault leaves
its own mark: a wide band (too long a step, a jumping cutoff, resonance),
a drift (loose constraints, resonance), a burst (wrapping, overlapping
atoms), a still band with the wrong kinetic energy (units), or a moving
centre of mass (drift kept).
\end{keyidea}

\notebook{11}{watch the healthy run, then break it in each of the eight ways}

\begin{selfcheck}
A run's kinetic energy at the start is 0.97\% of the value asked for, and
its energy is conserved well. What is the likely fault, and what factor
was missed?
\end{selfcheck}
